% !TeX root = ../proyecto.tex

\begin{UseCase}{CU-06}{Iniciar Sesión en el Sistema}{%
	Describe el proceso mediante el cual un usuario registrado autentica su identidad en el sistema web para acceder a las funcionalidades protegidas del Dashboard.
}
	\UCitem{Versión}{\color{Gray}
		1.0
	}\UCitem{Autor}{\color{Gray}
		Tesista de Ingeniería
	}\UCitem{Supervisa}{\color{Gray}
		Director de Trabajo Terminal
	}\UCitem{Actor}{
		Analista de Datos / Administrador del Sistema
	}\UCitem{Propósito}{\begin{Titemize}
		\Titem Verificar la identidad del usuario y establecer una sesión autenticada que habilite el acceso a todos los endpoints protegidos del sistema.
	\end{Titemize}
	}\UCitem{Entradas}{\begin{Titemize}
		\Titem Nombre de usuario (\texttt{username}).
		\Titem Contraseña (\texttt{password}).
		\Titem Opción ``Recordarme'' (booleano, opcional).
	\end{Titemize}
	}\UCitem{Origen}{\begin{Titemize}
		\Titem Formulario de inicio de sesión en la pantalla \IUref{IU-05}{Login}.
	\end{Titemize}
	}\UCitem{Salidas}{\begin{Titemize}
		\Titem Redirección al Dashboard principal \IUref{IU-01}{Dashboard} en caso de éxito.
		\Titem Mensaje de error visual si las credenciales son incorrectas o la cuenta está deshabilitada.
	\end{Titemize}
	}\UCitem{Destino}{\begin{Titemize}
		\Titem Sesión Flask-Login (cookie de sesión en el navegador).
		\Titem Tabla \texttt{users} en PostgreSQL (actualización de \texttt{last\_login}).
	\end{Titemize}
	}\UCitem{Precondiciones}{\begin{Titemize}
		\Titem El usuario debe estar previamente registrado en el sistema.
		\Titem La base de datos PostgreSQL debe estar operativa.
	\end{Titemize}
	}\UCitem{Postcondiciones}{\begin{Titemize}
		\Titem Se establece una sesión activa que permite el acceso a todas las rutas protegidas por \texttt{@login\_required}.
		\Titem Se registra la marca de tiempo de último inicio de sesión (\texttt{last\_login}).
	\end{Titemize}
	}\UCitem{Errores}{\begin{Titemize}
		\Titem {\bf \hypertarget{CU-06.E1}{E1}}: Credenciales incorrectas (usuario no encontrado o contraseña inválida) \Trayref{CU-06}{A}.
		\Titem {\bf \hypertarget{CU-06.E2}{E2}}: La cuenta del usuario está deshabilitada por el administrador \Trayref{CU-06}{B}.
	\end{Titemize}
	}\UCitem{Tipo}{
		Caso de uso primario
	}\UCitem{Observaciones}{
		La protección contra redirecciones abiertas (\textit{Open Redirect}) se implementa en el helper \texttt{\_is\_safe\_url()}, que valida que la URL de retorno pertenezca al mismo dominio del sistema antes de redirigir.
	}
\end{UseCase}

%--------------------------------------
\begin{UCtrayectoria}
	\UCpaso[\UCActor] Navega a la URL raíz del sistema; es redirigido automáticamente a la pantalla \IUref{IU-05}{Login}.
	\UCpaso[\UCActor] Introduce su nombre de usuario y contraseña en el formulario y hace clic en \IUbutton{Iniciar Sesión}.
	\UCpaso[\UCsist] Valida el formato del formulario (campos no vacíos).
	\UCpaso[\UCsist] Consulta la tabla \texttt{users} buscando el registro por \texttt{username} \ErrorRef{CU-06}{E1}{Credenciales inválidas}.
	\UCpaso[\UCsist] Verifica que la cuenta esté activa (\texttt{is\_active = True}) \ErrorRef{CU-06}{E2}{Cuenta deshabilitada}.
	\UCpaso[\UCsist] Valida el hash de la contraseña mediante \texttt{check\_password()}.
	\UCpaso[\UCsist] Establece la sesión de Flask-Login y actualiza el campo \texttt{last\_login}.
	\UCpaso[\UCsist] Redirige al Dashboard principal \IUref{IU-01}{Dashboard} o a la página solicitada originalmente (parámetro \texttt{next}).
\end{UCtrayectoria}

%--------------------------------------
\begin{UCtrayectoriaA}{CU-06}{A}{El nombre de usuario no existe o la contraseña es incorrecta}
	\UCpaso[\UCsist] Registra el intento fallido como \texttt{WARNING} en los logs del sistema, incluyendo la IP del cliente.
	\UCpaso[\UCsist] Muestra el mensaje ``Usuario o contraseña incorrectos'' en el formulario sin revelar cuál de los dos es incorrecto.
	\UCpaso[] El Caso de Uso termina; el actor puede reintentar.
\end{UCtrayectoriaA}

%--------------------------------------
\begin{UCtrayectoriaA}{CU-06}{B}{La cuenta del usuario está deshabilitada}
	\UCpaso[\UCsist] Detecta \texttt{is\_active = False} en el registro del usuario.
	\UCpaso[\UCsist] Muestra el mensaje ``Tu cuenta está deshabilitada. Contacta al administrador.''
	\UCpaso[\UCActor] Contacta al administrador para solicitar la reactivación.
	\UCpaso[] El Caso de Uso termina sin establecer sesión.
\end{UCtrayectoriaA}
